\section{Proven reference bounds}
The following results validate the formulations and are not claimed as new
contributions to graph classification. Labeling of Cartesian-product classes
has already been studied \cite{calamoneri}. Complete proofs are included
to avoid inferring formulas over infinite families from finite observations.

\begin{proposition}\label{prop:equal-tree}
For a nonempty tree $T$ and integer $h=k\geq1$,
$\lambda_{h,h}(T)=h\Delta(T)$.
\end{proposition}
\begin{proof}
The closed neighborhood of a degree-$\Delta$ vertex contains $\Delta+1$
vertices at pairwise distance at most 2, requiring span at least $h\Delta$.
For the upper bound, root the tree and color its square with $\Delta+1$
colors, level by level. Assign distinct colors to the children of each
vertex, excluding the colors of that vertex and its parent, if present.
A non-root vertex has at most $\Delta-1$ children and at least $\Delta-1$
available colors; the root has at most $\Delta$ children and $\Delta$
available colors. Distance-1 or distance-2 pairs are parent--child,
grandparent--grandchild, or sibling pairs, so this coloring is valid.
Multiplying colors $0,\ldots,\Delta$ by $h$ gives a labeling attaining
the bound. A one-vertex tree has span 0 and also satisfies the formula.
\end{proof}

For any graph with $n\geq1$ vertices, assigning distinct labels
$0,q,\ldots,(n-1)q$, where $q=\max(h,k)$, gives the elementary upper bound
$\lambda_{h,k}\leq(n-1)q$. This bound can be loose and differs from the
upper bound supplied by a specific heuristic labeling. In the comparisons,
the degree bound $h+(\Delta-1)\min(h,k)$ is also recorded separately
from final solver bounds.

\begin{proposition}\label{prop:extreme}
In an $L(2,1)$-labeling using labels $0,\ldots,\Delta+1$, every vertex
of degree $\Delta$ must receive label 0 or $\Delta+1$.
\end{proposition}
\begin{proof}
Its neighbors require $\Delta$ distinct labels because each pair is adjacent
or at distance 2. If the central label is $a\in\{1,\ldots,\Delta\}$,
labels $a-1,a,a+1$ are forbidden on its neighbors. Only $\Delta-1$ labels
remain, which is insufficient. The central label must therefore be a domain endpoint.
\end{proof}

\begin{proposition}\label{prop:grid}
For $n,m\geq4$, $\lambda_{2,1}(P_n\square P_m)=6$.
\end{proposition}
\begin{proof}
Index coordinates from 0 and define $f(i,j)=(2i+3j)\bmod7$, with
representatives in $\{0,\ldots,6\}$. Across an edge, the label difference
modulo 7 is $\pm2$ or $\pm3$, so the absolute difference between
representatives is at least 2. At distance 2, coordinate displacements are
$(\pm2,0),(0,\pm2),(\pm1,\pm1)$; the differences modulo 7 belong to
$\{\pm4,\pm6,\pm1,\pm5\}$ and are nonzero. This gives a valid labeling
of span at most 6.

Suppose a labeling has span at most 5 and translate it into $0,\ldots,5$.
The three internal vertices $(1,1),(1,2),(2,2)$ all have degree 4.
Proposition~\ref{prop:extreme} restricts them to labels 0 and 5. The two
successive edges force alternating labels, giving equal labels to the
endpoints. Their distance is 2, a contradiction. Thus the span is at least 6.
\end{proof}
The bound $\Delta+1=5$ alone is insufficient; the configuration of three
maximum-degree vertices is essential. This proof makes no equality claim
for $n<4$ or $m<4$, although the construction still provides a valid upper bound.

\begin{proposition}\label{prop:corona-lower}
For $n\geq3,m\geq1$, $\lambda_{2,1}(C_n\circ P_m)\geq m+4$.
\end{proposition}
\begin{proof}
Every core vertex has degree $\Delta=m+2$, giving the degree lower bound
$m+3$. If the span were at most $m+3$, Proposition~\ref{prop:extreme}
would restrict all core vertices to labels 0 and $m+3$. Three consecutive
core vertices are pairwise adjacent or at distance exactly 2 and therefore
require three distinct labels, a contradiction.
\end{proof}
This establishes only a lower bound. Observing $m+4$ over a finite sweep
does not replace an upper-bound construction for all $n,m$ and is
insufficient to establish general equality.
